\section{Tratamiento de errores léxicos}

\begin{longtable}{L{3.2cm}L{2.3cm}L{4.7cm}L{4.5cm}}
\caption{Errores, mensajes y estrategia de recuperación.}\label{tab:errores}\\
\toprule
\tableheader{Error} & \tableheader{Ejemplo} & \tableheader{Mensaje} & \tableheader{Recuperación} \\
\midrule
\endfirsthead
\toprule
\tableheader{Error} & \tableheader{Ejemplo} & \tableheader{Mensaje} & \tableheader{Recuperación} \\
\midrule
\endhead
Átomo citado sin cierre & \texttt{'juan} & \texttt{1:1 Atomo citado sin cierre} & Consume hasta salto de línea o EOF y continúa. \\
\rowcolor{palegray} Cadena sin cierre & \texttt{"hola} & \texttt{1:1 Cadena sin cierre} & Consume hasta salto de línea o EOF y continúa. \\
Comentario sin cierre & \verb|/* texto| & \texttt{1:1 Comentario de bloque sin cierre} & Consume hasta EOF; no existe un límite seguro posterior. \\
\rowcolor{palegray} Carácter no admitido & \texttt{@} & \texttt{1:1 Caracter no admitido} & Consume un carácter y continúa. \\
Número con sufijo & \texttt{12abc} & \texttt{1:1 Numero mal formado} & Consume el fragmento ASCII; conserva el punto que no preceda a un dígito. \\
\rowcolor{palegray} Exponente incompleto & \texttt{2e} & \texttt{1:1 Numero mal formado} & Consume el fragmento numérico completo. \\
Puntos repetidos & \texttt{3.14.15} & \texttt{1:1 Numero mal formado} & Consume el fragmento numérico completo. \\
\rowcolor{palegray} Escape no admitido & \verb|"x\q"| & \texttt{1:1 Secuencia de escape no admitida} & Consume hasta la comilla de cierre. \\
\bottomrule
\end{longtable}

\section{Plan de pruebas y resultados}

Los resultados se generaron el \FechaValidacion{} con Python \PythonValidacion{} mediante \path{scripts/validate.py}. La ejecución automatizada, reproducida en el anexo, completó \PruebasTotal{} pruebas con estado \texttt{OK}. Se conservan \PruebasValidasBase{} pruebas válidas y \PruebasInvalidasBase{} pruebas de entradas inválidas, posiciones o recuperación del corpus original; se añaden \PruebasRegresion{} pruebas de regresión del lexer y \PruebasCLI{} pruebas de la interfaz. El archivo válido produjo \TokensValidos{} tokens y \ErroresValidos{} errores; el archivo con errores produjo \TokensInvalidos{} tokens, \ErroresInvalidos{} diagnósticos y código de salida 1. Las cifras se incorporan desde \path{resultados.tex}, generado a partir de la ejecución, y el registro detallado queda en \path{docs/validation.json}.

\footnotesize
\begin{longtable}{C{0.9cm}L{2.6cm}L{2.75cm}L{3.0cm}L{3.0cm}C{1.3cm}}
\caption{Selección representativa del corpus automatizado.}\label{tab:pruebas}\\
\toprule
\tableheader{ID} & \tableheader{Entrada} & \tableheader{Objetivo} & \tableheader{Esperado} & \tableheader{Obtenido} & \tableheader{Estado} \\
\midrule
\endfirsthead
\toprule
\tableheader{ID} & \tableheader{Entrada} & \tableheader{Objetivo} & \tableheader{Esperado} & \tableheader{Obtenido} & \tableheader{Estado} \\
\midrule
\endhead
P01 & \texttt{padre} & Átomo simple & ATOMO & ATOMO & OK \\
\rowcolor{palegray} P03 & \texttt{'Juan Perez'} & Átomo citado & \makecell[l]{ATOMO\_\\CITADO} & \makecell[l]{ATOMO\_\\CITADO} & OK \\
P07 & \texttt{\_} & Anónima sin índice & \makecell[l]{VARIABLE\_\\ANONIMA} & \makecell[l]{VARIABLE\_\\ANONIMA} & OK \\
\rowcolor{palegray} P10 & \texttt{2e-3} & Real con exponente & REAL & REAL & OK \\
P16 & \verb|\==| & Máxima coincidencia & \verb|\==| & \verb|\==| & OK \\
\rowcolor{palegray} P24 & \texttt{is island} & Límite de palabra & OP / ATOMO & OP / ATOMO & OK \\
P27 & \texttt{padre(X)} repetido & Sin duplicados & Índices 0 y 0 & Índices 0 y 0 & OK \\
\rowcolor{palegray} P29 & \texttt{'juan} & Cierre faltante & Error & Error recuperado & OK \\
P31 & \verb|/* texto| & Comentario abierto & Error & Error hasta EOF & OK \\
\rowcolor{palegray} P33 & \texttt{12abc} & Número mal formado & Error & Error recuperado & OK \\
P35 & \texttt{3.14.15} & Dos puntos & Error & Error recuperado & OK \\
\rowcolor{palegray} P38 & \verb|a.\n  @| & Ubicación & 2:3 & 2:3 & OK \\
P39 & \texttt{1e+} & Exponente incompleto & Fragmento \texttt{1e+} & Fragmento \texttt{1e+} & OK \\
\rowcolor{palegray} P40 & \texttt{12abc, ok.} & Recuperación numérica & \texttt{, ok .} & \texttt{, ok .} & OK \\
\bottomrule
\end{longtable}
\normalsize

\subsection{Análisis de resultados}

El registro JSON incluye huellas SHA-256 de los archivos de código, pruebas, corpus y del propio validador. Antes de calcularlas se normalizan los finales CRLF a LF para comparar copias entre plataformas. El validador comprueba que estas entradas no hayan cambiado durante la ejecución. Las huellas permiten identificar el contenido probado y detectar evidencia desactualizada; no acreditan autoría ni sustituyen las contribuciones registradas en Git.

La cobertura comprueba cada familia léxica, los límites entre palabras y operadores, la deduplicación de lexemas y la actualización de líneas y columnas. Los casos P16 y P24 muestran que la prioridad no se resuelve solo por el orden de las alternativas: primero se consume el lexema completo y después se clasifica. Los casos P39 y P40 verifican que un exponente incompleto se reporte como un solo fragmento y que el analizador retome el reconocimiento en el delimitador siguiente.

El archivo con errores confirma una recuperación parcial. Después de \texttt{@}, \texttt{12abc}, una cadena abierta al final de línea, \texttt{3.14.15} y un escape inválido, el lexer vuelve a reconocer tokens posteriores. En cambio, un comentario de bloque abierto consume hasta el final del archivo porque no existe una marca léxica segura que indique dónde habría terminado.

Las regresiones verifican además saltos LF, CRLF y CR, preservación del punto de cláusula después de números mal formados y fragmentos de literales que terminan en una barra invertida. La interfaz se comprueba ejecutando procesos reales: formato solicitado en el enunciado, referencias a la tabla de lexemas, recuperación en el corpus inválido, BOM UTF-8, texto Unicode incluso con salida redirigida, archivo vacío, ayuda, archivo inexistente y codificación inválida. Un error de lectura genera un diagnóstico de uso y código 2 sin una traza interna de Python.

La principal limitación está en el alcance: el programa usa identificadores ASCII y no intenta reproducir toda la sintaxis de SWI-Prolog. Tampoco decide si una secuencia de tokens forma una cláusula válida. En consecuencia, los resultados demuestran conformidad con el subconjunto definido en este informe, no compatibilidad total con ISO Prolog.
